%(BEGIN_QUESTION)
% Copyright 2006, Tony R. Kuphaldt, released under the Creative Commons Attribution License (v 1.0)
% This means you may do almost anything with this work of mine, so long as you give me proper credit

Et svært nyttig prinsipp i fysikk er \textit{den ideelle gassloven}, kalt dette fordi den knytter trykk, volum, stoffmengde og temperatur for en ideell gass sammen i ett ryddig matematisk uttrykk:

$$PV = nRT$$

\noindent
Der,

$P$ = absolutt trykk (atmosfærer)

$V$ = volum (liter)

$n$ = gassmengde (mol)

$R$ = universell gasskonstant (0.0821 L $\cdot$ atm / mol $\cdot$ K)

$T$ = absolutt temperatur (K)

\vskip 10pt

Merk at temperaturen $T$ i denne likningen må være i \textit{absolutte} enheter (Kelvin). Omskriv den ideelle gassloven slik at den kan bruke en verdi for $T$ i enheten $^{o}$C.

Omskriv deretter likningen én gang til slik at den kan bruke en verdi for $T$ i enheten $^{o}$F.

\underbar{file i00342}
%(END_QUESTION)





%(BEGIN_ANSWER)

$$PV = nR(T + 273.15) \hbox{\hskip 30pt Temperatur i grader C}$$

$$PV = nR\left({5 \over 9}(T - 32) + 273.15\right) \hbox{\hskip 30pt Temperatur i grader F}$$

Det bør bemerkes at virkelige gasser avviker betydelig fra modellen for ideell gass ved temperaturer nær absolutt null, spesielt når faseendringer (kondensering og/eller størkning) skjer. Likevel er det mulig med fysikalsk intuisjon å se at temperaturen $T$ i denne likningen må være absolutt og ikke relativ, og det er kjernen i denne utfordringsoppgaven.

%(END_ANSWER)





%(BEGIN_NOTES)


%INDEX% Physics, heat and temperature: absolute temperature
%INDEX% Physics, static fluids: ideal gas law

%(END_NOTES)
